% 子文件：2.极限与连续.tex
% 在主文档使用 \input{2.极限与连续} 引入，不要写documentclass与document环境

\section{数列的极限}

\section{函数的极限}

\section{极限的性质，无穷小与无穷大}
\subsection{极限的性质}
定理2.3\quad 如果极限，必唯一

\subsection{无穷小与无穷大}
定理2.6\quad 有限个无穷小之和仍为无穷小

定理2.7\quad 无穷小与

\section{极限的运算法则}

\section{极限存在准则，两个重要极限}

\section{无穷小的比较}
等价无穷小\quad 牢记，当$x\to0$时，下列各对无穷小是等价的
\[
\begin{aligned}
\sin x &\sim x, & \tan x &\sim x, & \sqrt{1+x}-1 &\sim \frac12 x, & 1-\cos x &\sim \frac12 x^2,\\
\arctan x &\sim x, & \ln(1+x) &\sim x, & e^x-1 &\sim x, & \tan x-\sin x &\sim \frac12 x^3
\end{aligned}
\]

\section{函数的连续性}
\subsection{连续与间断}
定义 2.10\quad 设

\paragraph{第一类间断点}
左、右极限$f(x_0^-)$和$f(x_0^+)$都存在的间断点$x_0$，称为第一类间断点。
\begin{enumerate}
    \item $f(x_0^-)\neq f(x_0^+)$，即左、右极限都存在，但不相等，不管在$x_0$处函数是否有定义，这种第一类间断点叫做跳跃间断点。
    \item $f(x_0^-)＝f(x_0^+)$，但不等于$f(x_0)$或$f(x_0)$不存在，即有极限而不连续。这种第一类间断点叫做可去间断点。（只要补充或修改函数在$x_0$处的定义，令，就可以得到在$x_0$处连续的函数。）
\end{enumerate}

\paragraph{第二类间断点}
左、右极限至少有一个不存在的间断点，叫做第二类间断点。
\begin{itemize}
    \item 无穷间断点
    \item 振荡间断点
\end{itemize}

\subsection{函数连续性的判定定理}

\subsection{连续在极限运算中的应用}

\subsection{闭区间上连续函数的性质}
定理2.23

定义2.11

定理2.25